% Zone „Das kennst du schon“ – aus bank/flaechen/zone.jsonl (zusammenbau v0.1)
\einheitenkopf[zone][Kennst du schon]{Das kennst du schon}
\setcounter{aufgabe}{0}
%% TODO Grundvorstellungs-Aufgabe als erste Hauptnummer der Zone (2.8) fehlt in der Bank

%% TODO Titel: Ich-kann-Satz fehlt in der Bank (Platzhalter: Kettenname) – Nr. 1
%% TODO Anweisung: ein Satz über der ersten Teilaufgabe fehlt in der Bank – Nr. 1
\begin{aufgabe}{Längeneinheiten umrechnen (mm, cm, m), gemischte Angaben angleichen}
%% TODO Darstellung für schwach fehlt in der Bank (flaechen-zone-f1-v1; Abschnitt „Für schwache Schüler“)
\swz{$3$ m – wie viele cm?}{}
%% TODO Darstellung für schwach fehlt in der Bank (flaechen-zone-f1-v3; Abschnitt „Für schwache Schüler“)
\swz{$1{,}35$ m – wie viele cm?}{}
\end{aufgabe}

%% TODO Titel: Ich-kann-Satz fehlt in der Bank (Platzhalter: Kettenname) – Nr. 2
%% TODO Anweisung: ein Satz über der ersten Teilaufgabe fehlt in der Bank – Nr. 2
\begin{aufgabe}{Flächeneinheiten (cm², m²) und Umrechnungszahl 100}
%% TODO Darstellung für schwach fehlt in der Bank (flaechen-zone-f2-v1; Abschnitt „Für schwache Schüler“)
\swz{$4$ cm² – wie viele mm²?}{}
%% TODO Darstellung für schwach fehlt in der Bank (flaechen-zone-f2-v3; Abschnitt „Für schwache Schüler“)
\swz{$250$ cm² – wie viele dm²?}{}
\end{aufgabe}

%% TODO Titel: Ich-kann-Satz fehlt in der Bank (Platzhalter: Kettenname) – Nr. 3
%% TODO Anweisung: ein Satz über der ersten Teilaufgabe fehlt in der Bank – Nr. 3
\begin{aufgabe}{Multiplizieren und Dividieren mit Dezimalzahlen (Kommazahl mal Kommazahl, zweistellig geteilt durch einstellig)}
%% TODO Darstellung für schwach fehlt in der Bank (flaechen-zone-f3-v1; Abschnitt „Für schwache Schüler“)
\swz{$0{,}5 \cdot 8$ – wie viel?}{}
%% TODO Darstellung für schwach fehlt in der Bank (flaechen-zone-f3-v3; Abschnitt „Für schwache Schüler“)
\swz{$2{,}5 \cdot 1{,}4$ – wie viel?}{}
\end{aufgabe}

%% TODO Titel: Ich-kann-Satz fehlt in der Bank (Platzhalter: Kettenname) – Nr. 4
%% TODO Anweisung: ein Satz über der ersten Teilaufgabe fehlt in der Bank – Nr. 4
\begin{aufgabe}{Formel nach einer Größe umstellen (A = a · b → b = A : a)}
%% TODO Darstellung für schwach fehlt in der Bank (flaechen-zone-f4-v1; Abschnitt „Für schwache Schüler“)
\swz{$P = q \cdot r$ – nach $r$ umstellen?}{}
%% TODO Darstellung für schwach fehlt in der Bank (flaechen-zone-f4-v3; Abschnitt „Für schwache Schüler“)
\swz{$G = m \cdot p$ mit $G = 7{,}2$ und $m = 3$ – $p$?}{}
\end{aufgabe}

%% TODO Titel: Ich-kann-Satz fehlt in der Bank (Platzhalter: Kettenname) – Nr. 5
%% TODO Anweisung: ein Satz über der ersten Teilaufgabe fehlt in der Bank – Nr. 5
\begin{aufgabe}{Wurzel ziehen bei Quadratzahlen (√49)}
%% TODO Darstellung für schwach fehlt in der Bank (flaechen-zone-f5-v1; Abschnitt „Für schwache Schüler“)
\swz{$\sqrt{25}$ – wie viel?}{}
%% TODO Darstellung für schwach fehlt in der Bank (flaechen-zone-f5-v3; Abschnitt „Für schwache Schüler“)
\swz{$\sqrt{144}$ – wie viel?}{}
\end{aufgabe}

%% TODO Titel: Ich-kann-Satz fehlt in der Bank (Platzhalter: Kettenname) – Nr. 6
%% TODO Anweisung: ein Satz über der ersten Teilaufgabe fehlt in der Bank – Nr. 6
\begin{aufgabe}{Rechten Winkel erkennen und einzeichnen (Geodreieck)}
\begin{teile}
\teil Ist das ein rechter Winkel? \janein
\end{teile}
\swa{An welcher Ecke ist ein rechter Winkel? Prüfe mit dem Geodreieck und markiere ihn.}{\viereck{(0,0)}{(5,0)}{(5,3)}{(1,2.5)}}
\end{aufgabe}

%% TODO Titel: Ich-kann-Satz fehlt in der Bank (Platzhalter: Kettenname) – Nr. 7
%% TODO Anweisung: ein Satz über der ersten Teilaufgabe fehlt in der Bank – Nr. 7
\begin{aufgabe}{Kreisfläche (π · r²)}
%% TODO Darstellung für schwach fehlt in der Bank (flaechen-zone-f7-v1; Abschnitt „Für schwache Schüler“)
\swz{Kreis mit Radius $r = 1$ m – Fläche? Runde auf zwei Stellen.}{}
%% TODO Darstellung für schwach fehlt in der Bank (flaechen-zone-f7-v3; Abschnitt „Für schwache Schüler“)
\swz{Kreis mit Radius $r = 1{,}5$ m – Fläche? Runde auf zwei Stellen.}{}
\end{aufgabe}

%% TODO Titel: Ich-kann-Satz fehlt in der Bank (Platzhalter: Kettenname) – Nr. 8
%% TODO Anweisung: ein Satz über der ersten Teilaufgabe fehlt in der Bank – Nr. 8
\begin{aufgabe}{Flächeneinheiten (cm², m²) und Umrechnungszahl 100 – Fehler finden}
\swz[3]{Mia rechnet: $8$ cm² $= 80$ mm². Wo ist der Fehler?}{}
\end{aufgabe}

%% TODO Titel: Ich-kann-Satz fehlt in der Bank (Platzhalter: Kettenname) – Nr. 9
%% TODO Anweisung: ein Satz über der ersten Teilaufgabe fehlt in der Bank – Nr. 9
\begin{aufgabe}{Flächeneinheiten (cm², m²) und Umrechnungszahl 100 – selbst rechnen}
%% TODO Darstellung für schwach fehlt in der Bank (flaechen-zone-f2-v6; Abschnitt „Für schwache Schüler“)
\swz{$9$ dm² – wie viele cm²?}{}
\end{aufgabe}
